\section{Related Work}
\label{sec:related}

\textbf{On-policy distillation.}
On-policy distillation trains a student on its own generations while querying a teacher for dense token-level supervision \citep{agarwal2024policy}. Recent work connects this objective to policy-gradient learning \citep{vopd2026,exopd2026}, organizes distillation methods by their rollout source and divergence direction \citep{decoupled2026}, and studies how the choice of divergence affects policy improvement \citep{distil2026}. Other analyses show that the effectiveness of OPD depends strongly on student--teacher support overlap and that teacher supervision can become unreliable on student-generated prefixes \citep{manyfaces2026,rethinking2026}. Several methods address this issue at the token level by adapting the KL direction \citep{aopd2026,hpd2026}, selecting supported tokens \citep{taopd2026}, or mixing teacher and student prefixes \citep{xu2025skd,brts2026}. Orthogonal to where supervision is applied, a further group of variants changes what the teacher provides: rubrics or preference pairs induced from teacher text stand in for logits when the teacher is black-box \citep{ropd2026,orpodistill2025}, privileged context conditions the teacher's scores \citep{mopd2026}, and another line couples teacher supervision with an explicit task reward inside a single objective \citep{kdrl2025,rlad2026,sparse2dense2026,chord2026}, following earlier on-policy policy distillation for control \citep{ppd2025}. Teacher uncertainty is also used to adapt token-level distillation objectives \citep{eopd2026,egrsd2026}, while other work obtains corrective supervision by probing alternative continuations or backtracking to an earlier reasoning prefix \citep{spot2026,motab2026}. UOPD selects corrections during environment interaction, where executing a corrected action also changes the states visited next.

\textbf{Multi-turn agent distillation and intervention.}
In multi-turn environments, errors alter subsequent observations and cause distribution shift to accumulate across turns \citep{wang2026exploring,sod2026,sageopd2026}. Existing approaches mitigate this drift by reweighting or filtering supervision \citep{sod2026,pruneopd2026,sdar2026}, limiting or adapting rollout exposure \citep{wang2026exploring,fullrollouts2026,adwin2026,turnopd2026}, or modifying the roll-in trajectory through teacher prefixes, refinement, or lookahead feedback \citep{reopd2026,trd2026,lgr2026,serl2026}. SAGE-OPD \citep{sageopd2026} uses teacher judgments and confidence to scale turn-level supervision. Teacher actions also guide rollouts through scheduled turn mixing \citep{guidedopd2026} or corrections selected and validated after rollout collection \citep{futurebridgeopd2026}. UOPD uses teacher uncertainty to target critical steps and intervene before the student's proposed action is executed, providing a corrective learning target and redirecting the subsequent rollout. UOPD also connects to DAgger-style expert supervision at learner-visited states \citep{ross2011dagger,distil2026} and HG-DAgger's expert takeover \citep{kelly2019hgdagger}. It automates takeover using teacher uncertainty on student actions and trains the student on the executed corrections.
